% GENERATED FILE. Edit canonical lessons, then run npm run generate:books.
% canonical-source-hash: fnv1a64:b9bdbcc5be1ca552
% canonical-lessons: MR-C324-nidhan, MR-C324-sraddhanjali, MR-C324-ayojan-karne, MR-C324-par-padne, MR-C324-racne

\chapter{Remembering a Life}
\label{ch:mr-remembering-a-life}
\hlchaptermodality{324}

\begin{chapteropening}
\textbf{By the end of this chapter:} \emph{I can talk about a death, a tribute and the gathering that follows.}

Complete the last lesson of chapter 324: I can talk about a death, a tribute and the gathering that follows.
\end{chapteropening}

\section[\texorpdfstring{\mr{निधन} (nidhan)}{निधन (nidhan)}]{\mr{निधन} (nidhan) — a passing; a death}

\label{lesson:MR-C324-nidhan}



\begin{quote}
Before the new one: say the Marathi for an anniversary, then the Marathi for a closing.
\end{quote}

\subsection*{You'll want to know: \mr{निधन}}
\textbf{\mr{निधन}} — \emph{nidhan} — \textquotedblleft{}a passing; a death\textquotedblright{}.

\begin{quote}
\mr{प्रसिद्ध} \mr{गायकाच्या} \mr{निधनाची} \mr{बातमी} \mr{काल} \mr{आली}.
\end{quote}

\begin{quote}
\emph{prasiddha gāykācyā nidhanācī bātmī kāl ālī.}
\end{quote}

\begin{quote}
News of the famous singer's passing came yesterday.
\end{quote}

Neuter.

The respectful word: X-\mr{चं} \mr{निधन} \mr{झालं}.

\begin{cousinweb}
\textbf{\mr{निधन}} (\emph{nidhan}) is a tatsama, taken from Sanskrit \emph{nidhana}, 'a settling, a close; death' (Monier-Williams).
\end{cousinweb}

\begin{grammarlens}[title={What you've built}]
One more word for telling what happened.
\end{grammarlens}

\subsection*{Your turn}
\begin{itemize}
\raggedright
  \item \emph{Say it:} \emph{nidhan}
  \item \emph{Say it:} \emph{nidhan}, once more
  \item \emph{Say it:} \emph{prasiddha gāykācyā nidhanācī bātmī kāl ālī.} — News of the famous singer's passing came yesterday.
  \item \emph{Recall:} say the Marathi for an anniversary, then the Marathi for a closing, then say \emph{nidhan} again
\end{itemize}

\begin{tcolorbox}[breakable,colback=teal!4,colframe=teal!35!black,title={Before you move on}]
What does \mr{निधन} mean? (\textquotedblleft{}A passing; a death\textquotedblright{}.) Say it once more.
\end{tcolorbox}
\section[\texorpdfstring{\mr{श्रद्धांजली} (śraddhāñjalī)}{श्रद्धांजली (śraddhāñjalī)}]{\mr{श्रद्धांजली} (śraddhāñjalī) — a tribute}

\label{lesson:MR-C324-sraddhanjali}



\begin{quote}
Before the new one: say the Marathi for a closing, then the Marathi for a passing.
\end{quote}

\subsection*{You'll want to know: \mr{श्रद्धांजली}}
\textbf{\mr{श्रद्धांजली}} — \emph{śraddhāñjalī} — \textquotedblleft{}a tribute\textquotedblright{}.

\begin{quote}
\mr{सभेत} \mr{सगळ्यांनी} \mr{त्यांना} \mr{श्रद्धांजली} \mr{वाहिली}.
\end{quote}

\begin{quote}
\emph{sabhet sagaḷyānnī tyānnā śraddhāñjalī vāhilī.}
\end{quote}

\begin{quote}
At the meeting everyone paid tribute to him.
\end{quote}

Feminine.

\mr{श्रद्धांजली} \mr{वाहणे}, to pay tribute.

\begin{cousinweb}
\textbf{\mr{श्रद्धांजली}} (\emph{śraddhāñjalī}) is built from Sanskrit \emph{śraddhā}, 'faith, reverence', and \emph{añjali}, 'hands held together in offering' (Monier-Williams).
\end{cousinweb}

\begin{grammarlens}[title={What you've built}]
One more word for telling what happened.
\end{grammarlens}

\subsection*{Your turn}
\begin{itemize}
\raggedright
  \item \emph{Say it:} \emph{śraddhāñjalī}
  \item \emph{Say it:} \emph{śraddhāñjalī}, once more
  \item \emph{Say it:} \emph{sabhet sagaḷyānnī tyānnā śraddhāñjalī vāhilī.} — At the meeting everyone paid tribute to him.
  \item \emph{Recall:} say the Marathi for a closing, then the Marathi for a passing, then say \emph{śraddhāñjalī} again
\end{itemize}

\begin{tcolorbox}[breakable,colback=teal!4,colframe=teal!35!black,title={Before you move on}]
What does \mr{श्रद्धांजली} mean? (\textquotedblleft{}A tribute\textquotedblright{}.) Say it once more.
\end{tcolorbox}
\section[\texorpdfstring{\mr{आयोजन} \mr{करणे} (āyojan karṇe)}{आयोजन करणे (āyojan karṇe)}]{\mr{आयोजन} \mr{करणे} (āyojan karṇe) — to organise}

\label{lesson:MR-C324-ayojan-karne}



\begin{quote}
Before the new one: say the Marathi for a passing, then the Marathi for a tribute.
\end{quote}

\subsection*{You'll want to know: \mr{आयोजन} \mr{करणे}}
\textbf{\mr{आयोजन} \mr{करणे}} — \emph{āyojan karṇe} — \textquotedblleft{}to organise\textquotedblright{}.

\begin{quote}
\mr{मित्रांनी} \mr{त्यांच्या} \mr{आठवणीत} \mr{एका} \mr{सभेचं} \mr{आयोजन} \mr{केलं}.
\end{quote}

\begin{quote}
\emph{mitrānnī tyāñcyā āṭhavṇīt ekā sabhecaṁ āyojan kelaṁ.}
\end{quote}

\begin{quote}
Friends organised a meeting in his memory.
\end{quote}

Past, with -\mr{ने} (he …): \mr{त्याने} \mr{आयोजन} \mr{केलं}, \emph{tyāne āyojan kelaṁ}.

What is organised takes -\mr{चं}: \mr{स्पर्धेचं} \mr{आयोजन} \mr{करणे}.

\begin{cousinweb}
\textbf{\mr{आयोजन} \mr{करणे}} (\emph{āyojan karṇe}) joins \emph{āyojan}, a tatsama taken from Sanskrit \emph{āyojana}, 'a joining, a bringing together' (Monier-Williams), to \emph{karṇe}, 'to do'.
\end{cousinweb}

\begin{grammarlens}[title={What you've built}]
One more word for telling what happened.
\end{grammarlens}

\subsection*{Your turn}
\begin{itemize}
\raggedright
  \item \emph{Say it:} \emph{āyojan karṇe}
  \item \emph{Say it:} \emph{āyojan karṇe}, once more
  \item \emph{Say it:} \emph{mitrānnī tyāñcyā āṭhavṇīt ekā sabhecaṁ āyojan kelaṁ.} — Friends organised a meeting in his memory.
  \item \emph{Recall:} say the Marathi for a passing, then the Marathi for a tribute, then say \emph{āyojan karṇe} again
\end{itemize}

\begin{tcolorbox}[breakable,colback=teal!4,colframe=teal!35!black,title={Before you move on}]
What does \mr{आयोजन} \mr{करणे} mean? (\textquotedblleft{}To organise\textquotedblright{}.) Say it once more.
\end{tcolorbox}
\section[\texorpdfstring{\mr{पार} \mr{पडणे} (pār paḍṇe)}{पार पडणे (pār paḍṇe)}]{\mr{पार} \mr{पडणे} (pār paḍṇe) — to go off; to be carried through}

\label{lesson:MR-C324-par-padne}



\begin{quote}
Before the new one: say the Marathi for a tribute, then the Marathi for to organise.
\end{quote}

\subsection*{You'll want to know: \mr{पार} \mr{पडणे}}
\textbf{\mr{पार} \mr{पडणे}} — \emph{pār paḍṇe} — \textquotedblleft{}to go off; to be carried through\textquotedblright{}.

\begin{quote}
\mr{सगळा} \mr{कार्यक्रम} \mr{शांतपणे} \mr{पार} \mr{पडला}.
\end{quote}

\begin{quote}
\emph{sagḷā kāryakram śāntapṇe pār paḍlā.}
\end{quote}

\begin{quote}
The whole event went off peacefully.
\end{quote}

Past (he …): \mr{तो} \mr{पार} \mr{पडला}, \emph{to pār paḍlā}.

Intransitive, said of events: \mr{लग्न} \mr{छान} \mr{पार} \mr{पडलं}.

\begin{cousinweb}
\textbf{\mr{पार} \mr{पडणे}} (\emph{pār paḍṇe}) is built from \emph{pār}, 'the far bank, the end', inherited from Sanskrit \emph{pāra} (CDIAL 8100), and \emph{paḍṇe}, 'to fall': to come through to the other side.
\end{cousinweb}

\begin{grammarlens}[title={What you've built}]
One more word for telling what happened.
\end{grammarlens}

\subsection*{Your turn}
\begin{itemize}
\raggedright
  \item \emph{Say it:} \emph{pār paḍṇe}
  \item \emph{Say it:} \emph{pār paḍṇe}, once more
  \item \emph{Say it:} \emph{sagḷā kāryakram śāntapṇe pār paḍlā.} — The whole event went off peacefully.
  \item \emph{Recall:} say the Marathi for a tribute, then the Marathi for to organise, then say \emph{pār paḍṇe} again
\end{itemize}

\begin{tcolorbox}[breakable,colback=teal!4,colframe=teal!35!black,title={Before you move on}]
What does \mr{पार} \mr{पडणे} mean? (\textquotedblleft{}To go off; to be carried through\textquotedblright{}.) Say it once more.
\end{tcolorbox}
\section[\texorpdfstring{\mr{रचणे} (racṇe)}{रचणे (racṇe)}]{\mr{रचणे} (racṇe) — to compose; to arrange}

\label{lesson:MR-C324-racne}



\begin{quote}
Before the new one: say the Marathi for to organise, then the Marathi for to go off.
\end{quote}

\subsection*{You'll want to know: \mr{रचणे}}
\textbf{\mr{रचणे}} — \emph{racṇe} — \textquotedblleft{}to compose; to arrange\textquotedblright{}.

\begin{quote}
\mr{त्यांच्या} \mr{मित्राने} \mr{त्यांच्यावर} \mr{एक} \mr{कविता} \mr{रचली}.
\end{quote}

\begin{quote}
\emph{tyāñcyā mitrāne tyāñcyāvar ek kavitā raclī.}
\end{quote}

\begin{quote}
His friend composed a poem about him.
\end{quote}

Past, with -\mr{ने} (he …): \mr{त्याने} \mr{रचलं}, \emph{tyāne raclaṁ}.

Transitive: \mr{कविता} \mr{रचणे}, to compose a poem; \mr{दगड} \mr{रचणे}, to stack stones.

\begin{cousinweb}
\textbf{\mr{रचणे}} (\emph{racṇe}) is built on Sanskrit \emph{racayati}, 'produces, arranges, composes' (Monier-Williams).
\end{cousinweb}

\begin{grammarlens}[title={What you've built}]
One more word for telling what happened.
\end{grammarlens}

\subsection*{Your turn}
\begin{itemize}
\raggedright
  \item \emph{Say it:} \emph{racṇe}
  \item \emph{Say it:} \emph{racṇe}, once more
  \item \emph{Say it:} \emph{tyāñcyā mitrāne tyāñcyāvar ek kavitā raclī.} — His friend composed a poem about him.
  \item \emph{Recall:} say the Marathi for to organise, then the Marathi for to go off, then say \emph{racṇe} again
\end{itemize}

\begin{tcolorbox}[breakable,colback=teal!4,colframe=teal!35!black,title={Before you move on}]
What does \mr{रचणे} mean? (\textquotedblleft{}To compose; to arrange\textquotedblright{}.) Say it once more.
\end{tcolorbox}
